\section{Simulation evidence: mathematical claims and counterexamples}
\label{sec:annex-p12-simulation-evidence}

\begin{annexitem}{annex:p12.simulation.error-chain}{thm:p12-simulation-error-chain}{The conditional error-chain bound}
\begin{proof}
Insert the exact model, supplied-input model and discretized model
between \(Y\) and \(\widehat m_h\). The difference is the sum of
the four differences named in the definition. The triangle inequality
bounds its magnitude by the sum of their four assumed bounds.
This algebra neither derives nor estimates any of those premises.
The previously proved numerical sign/order certificates apply to
the resulting enclosure of the same target quantity.
\end{proof}
\end{annexitem}

\begin{annexitem}{annex:p12.simulation.reference-comparison}{thm:p12-simulation-reference-comparison}{Upper and lower referent-comparison bounds}
\begin{proof}
Write
\(m-Y=(m-\widehat m)+(\widehat m-r)+(r-Y)\).
The triangle inequality gives the upper bound. Conversely,
\(\widehat m-r=(\widehat m-m)+(m-Y)+(Y-r)\), so
\(|\widehat m-r|\leq E_n+|m-Y|+E_r\).
Rearrange and combine with \(|m-Y|\geq0\) for the lower bound.
Comparing these endpoints with \(\tau\) proves both tolerance
conclusions, including their strict/non-strict boundary distinctions.
\end{proof}
\end{annexitem}

\begin{annexitem}{annex:p12.simulation.decay-euler}{thm:p12-simulation-decay-euler}{The exact model, Euler cases and a finite error bound}
\begin{proof}
Differentiate \(q_0\exp(-\lambda t)\) to verify the equation and
initial condition. For any other differentiable solution, the
derivative of \(\exp(\lambda t)q(t)\) is zero; its value is
therefore the constant \(q_0\). This proves uniqueness.
Induction on the recurrence gives \(v_N=q_0(1-a)^N\), where
\(a=\lambda h\). If \(0\leq a\leq1\), its multiplier lies
in \([0,1]\); the nonnegative monotonicity follows. For
\(1<a<2\), the multiplier lies in \((-1,0)\); at two it is
minus one, and above two it has magnitude greater than one.
These facts give each sign and magnitude claim.

For \(0<a<1\), define
\[
 r=-\ln(1-a)-a=\int_0^a\frac{t}{1-t}\,dt.
\]
The integral is nonnegative and at most
\(a^2/(2(1-a))\), since \((1-t)^{-1}\leq(1-a)^{-1}\)
on its integration interval. Thus
\[
 (1-a)^N=\exp(-Na)\exp(-Nr),\qquad
 0\leq1-\exp(-Nr)\leq Nr.
\]
For the last inequality, integrate \(\exp(-s)\leq1\) from
zero to \(Nr\). Multiplication by \(q(T)=q_0\exp(-Na)\)
gives an error at most \(q(T)Na^2/(2(1-a))\), which equals
the asserted expression because \(Na=\lambda T\).
When \(\lambda=0\), both formulas are identically \(q_0\).
\end{proof}
\end{annexitem}

\begin{annexitem}{annex:p12.simulation.arithmetic-propagation}{thm:p12-simulation-arithmetic-propagation}{The exact perturbed recurrence and its bound}
\begin{proof}
Set \(e_j=\widehat v_j-v_j\). Subtract the exact Euler recurrence
from the perturbed recurrence to obtain
\(e_{j+1}=(1-a)e_j+\delta_j\). Induction gives
\[
 e_N=(1-a)^Ne_0+\sum_{j=0}^{N-1}(1-a)^{N-1-j}\delta_j.
\]
Since \(1-a\in[0,1]\), the triangle inequality gives the stated
bound with nonnegative weights. For a constant upper bound \(\rho\)
and \(a>0\), summing the finite geometric series gives
\(\sum_{i=0}^{N-1}(1-a)^i=[1-(1-a)^N]/a\).
For \(a=0\), every coefficient equals one and the bound is
\(E_0+N\rho\). No distributional assumption was used.
\end{proof}
\end{annexitem}

\begin{annexitem}{annex:p12.simulation.decay-inputs}{thm:p12-simulation-decay-inputs}{A global bound on simultaneous parameter changes}
\begin{proof}
Change \(q_0\), then \(\lambda\), then \(t\), inserting the
two intermediate values in the difference. The first change has
magnitude at most \(|q_0-\widetilde q_0|\), since
\(\exp(-\lambda t)\leq1\). The derivative of
\(\widetilde q_0\exp(-\lambda t)\) with respect to
\(\lambda\) has magnitude
\(\widetilde q_0t\exp(-\lambda t)\leq QT\).
The derivative with respect to time, with the new parameter fixed,
has magnitude at most \(Q\Lambda\). Apply the one-variable
mean value theorem to these last two changes and add the three
bounds. The rectangle assumptions ensure these derivative bounds
hold along the entire two-segment parameter path.
\end{proof}
\end{annexitem}

\begin{annexitem}{annex:p12.simulation.boxbod-residual}{thm:p12-simulation-boxbod-residual}{A reference discrepancy established without rounding the exponential}
\begin{proof}
The stated decimal parameters satisfy \(\widehat B>213.8\)
and \(3\widehat\lambda>1.62\). For positive arguments the
exponential series has positive terms, and the exact rational sum
\[
 \sum_{j=0}^5\frac{(81/50)^j}{j!}
 =\frac{62759338217}{12500000000}>5.
\]
Hence \(\exp(3\widehat\lambda)>5\), so
\(1-\exp(-3\widehat\lambda)>4/5\).
The model prediction exceeds \(213.8(4/5)=171.04\).
Subtracting the recorded response \(149\) proves a discrepancy
greater than \(22.04\), and therefore greater than ten.
The argument compares the explicit model expression with the
explicit record. It introduces no premise about the unknown
experimental error or the actual generating mechanism.
\end{proof}
\end{annexitem}

\begin{annexitem}{annex:p12.simulation.finite-comparisons}{thm:p12-simulation-finite-comparisons}{Alternatives indistinguishable on a finite comparison set}
\begin{proof}
Let \(P(t)=\prod_{i=1}^n(t-t_i)^2\). It is a polynomial
vanishing at every \(t_i\), and \(P(s)>0\).
For any desired positive discrepancy \(A\), set
\(g(t)=f(t)+A P(t)/P(s)\).
Then \(g(t_i)=f(t_i)\) for all comparisons while
\(g(s)-f(s)=A\). A polynomial is continuous, smooth and analytic,
so addition preserves whichever stated regularity \(f\) has.
No position of \(s\) relative to the smallest and largest
comparison inputs was used, only its distinctness from them.
\end{proof}
\end{annexitem}

\begin{annexitem}{annex:p12.simulation.regime-change}{thm:p12-simulation-regime-change}{Indistinguishable earlier trajectories and different later decisions}
\begin{proof}
The two pieces defining \(q_2\) have the same value
\(q_0\exp(-\lambda T)\) at \(T\). Both are positive and
decreasing on their intervals, and their join is continuous and
nonincreasing. The same properties hold for \(q_1\).
For \(t>T\), their ratio is
\[
 q_2(t)/q_1(t)=\exp\bigl(- (\mu-\lambda)(t-T)\bigr),
\]
which is not one because \(\mu\ne\lambda\) and \(t-T>0\).
They agree for every earlier time by definition. Any evidence
restricted to that earlier trajectory therefore has exactly the
same values under both alternatives; it cannot distinguish them
without another premise. For the displayed numerical example,
\(e^{-2}>0.1>e^{-3}\) follows from \(e^2<10<e^3\).
For an exact check, the exponential series gives \(e>5/2\).
Also \(j!\geq2^{j-1}\) for \(j\geq2\), strictly for
\(j\geq3\); therefore
\(e=2+\sum_{j=2}^\infty1/j!<2+\sum_{j=2}^\infty2^{1-j}=3\).
It follows that \(e^2<9<10\) and \(e^3>125/8>10\).
\end{proof}
\end{annexitem}

\begin{annexitem}{annex:p12.simulation.ordinal-scores}{thm:p12-simulation-ordinal-scores}{A reversal under increasing relabeling}
\begin{proof}
Take profiles \(A=(0,4)\) and \(B=(2,3)\), whose label means
are two and \(2.5\). Replace the five ordered labels
\((0,1,2,3,4)\) by \((0,1,2,3,100)\), a strictly increasing
relabeling preserving every ordinal comparison. Their new means
are fifty and \(2.5\), reversing the ranking. Thus the ordering
of levels alone cannot determine the mean-based aggregate.
Choosing a numerical scale or aggregation rule requires an
additional interpretation; it is not supplied by the ordinal order.
\end{proof}
\end{annexitem}
